% Prova fictícia no formato da classe exam (dados de teste)
\documentclass[addpoints]{exam}
\usepackage{amsmath,amssymb}
\newcommand{\R}{\mathbb{R}}
\DeclareMathOperator{\sen}{sen}
\newcommand{\abs}[1]{\left|#1\right|}
\printanswers
\begin{document}
\begin{center}\textbf{Prova fictícia}\end{center}
\begin{questions}

\question[2] Quanto vale $\displaystyle\lim_{x\to 0}\frac{\sen x}{x}$?
\begin{choices}
  \choice $0$
  \CorrectChoice $1$
  \choice $\infty$
  \choice Não existe.
\end{choices}
\begin{solution}
Limite fundamental trigonométrico.
\end{solution}

\question[3] Considere $f\colon\R\to\R$, $f(x)=x^2$.
\begin{parts}
\part Calcule $f'(x)$.
\part Esboce o gráfico.
\end{parts}
\begin{solutionorbox}[2in]
$f'(x) = 2x$; parábola com vértice na origem.
\end{solutionorbox}

% \question Esta questão está comentada e não deve ser importada.

\titledquestion{Bônus}[1] O valor de $2+2$ é \fillin[$4$][2cm].

\question Assinale as corretas.
\begin{oneparchoices}
\CorrectChoice $\abs{-1}=1$ \choice $\abs{-1}=-1$
\end{oneparchoices}
\vspace{2cm}

\end{questions}
\end{document}
